\section{Check for a Balanced Binary Tree}
\label{sec:trees-balanced}

\problemstatement{A binary tree is height-balanced if, for \emph{every}
node, the heights of its left and right subtrees differ by at most $1$.
Given the root, determine whether the tree is balanced.}

\begin{patternbox}
The naive reading -- ``for every node, compute both subtree heights and
compare'' -- is an $O(n^2)$ trap on a skewed tree, because height is
recomputed from scratch at every node. The keyword \emph{``for every
node''} combined with \emph{``heights''} says: compute height once,
bottom-up, and \textbf{fold the balance check into that same pass}.
\end{patternbox}

\subsection*{Understanding the problem carefully}

Height and the balance condition are computed from the same recursion. If
each recursive call returns the height of its subtree, the parent can, in
$O(1)$, check whether its two child heights differ by more than $1$ before
returning its own height. A single sentinel value (say $-1$) propagated
upward can signal ``a violation was found below'' so the whole recursion
short-circuits.

\begin{ideabox}[Return Height, or $-1$ to Mean ``Already Unbalanced'']
Write a helper that returns the subtree height normally, but returns $-1$
the moment it detects imbalance -- either because a child already returned
$-1$, or because the two child heights differ by more than $1$. The final
answer is ``balanced'' iff the root's call returns something other than
$-1$. One pass, $O(n)$.
\end{ideabox}

\subsection*{Algorithm}

\begin{enumerate}
  \item \textsc{Check}(\textit{node}):
  \begin{enumerate}
    \item If \textit{node} is null, return $0$.
    \item $L = $ \textsc{Check}(\textit{node}.left); if $L=-1$ return
          $-1$.
    \item $R = $ \textsc{Check}(\textit{node}.right); if $R=-1$ return
          $-1$.
    \item If $|L-R|>1$ return $-1$; else return $1+\max(L,R)$.
  \end{enumerate}
  \item The tree is balanced iff \textsc{Check}(root)$\ne -1$.
\end{enumerate}

\subsection*{Dry run}

\dryrun{Tree: root $1$; left $2$ with a left child $3$ that has a left
child $4$; right subtree of $1$ empty.}

\begin{itemize}
  \item Height($4$)$=1$, Height($3$)$=1+\max(1,0)=2$.
  \item At $2$: left height $2$, right height $0$; $|2-0|=2>1
        \Rightarrow$ return $-1$.
  \item $-1$ propagates to $1$; answer: \textbf{not balanced}.
\end{itemize}

\subsection*{C++ implementation}

\begin{lstlisting}
#include <bits/stdc++.h>
using namespace std;

struct TreeNode {
    int val;
    TreeNode* left;
    TreeNode* right;
    TreeNode(int v) : val(v), left(nullptr), right(nullptr) {}
};

int check(TreeNode* node) {
    if (node == nullptr) return 0;
    int L = check(node->left);
    if (L == -1) return -1;
    int R = check(node->right);
    if (R == -1) return -1;
    if (abs(L - R) > 1) return -1;
    return 1 + max(L, R);
}

bool isBalanced(TreeNode* root) {
    return check(root) != -1;
}

int main() {
    TreeNode* root = new TreeNode(1);
    root->left = new TreeNode(2);
    root->left->left = new TreeNode(3);
    root->left->left->left = new TreeNode(4);

    cout << (isBalanced(root) ? "balanced" : "not balanced") << "\n";
    return 0;
}
\end{lstlisting}

\begin{complexitybox}
\textbf{Time Complexity.} $O(n)$ -- height and the balance test share one
postorder pass, versus $O(n^2)$ if height were recomputed per node.
\textbf{Space Complexity.} $O(h)$ for the recursion stack.
\end{complexitybox}

\begin{takeawaybox}
Overloading a return value with a sentinel ($-1$ = ``violation found'') is
a clean way to combine ``compute a quantity'' with ``check a property''
in one traversal, avoiding the quadratic blow-up of recomputing the
quantity at every node.
\end{takeawaybox}
